% Chapter 5

\chapter{Discussion}\label{Chapter5}

%----------------------------------------------------------------------------------------
%	SECTION 1
%----------------------------------------------------------------------------------------

\section{A Cognitive Severity Gradient, Not Discrete Phenotypes}
Overall, I found no evidence of discrete cognitive phenotypes in this cohort of people with acquired brain injury (ABI); instead, I found a clear cognitive-severity gradient. I identified patients as being in the `above average', `near normal' and `global impairment' tiers based upon their performance across the battery. Importantly, these tiers differ by level, not by profile shape. As illustrated in Figure~\ref{fig:radar_profiles}, each of the six domains moves in the same direction. One principal component explained 56\% of variance in the domain scores and loaded positively across all six domains with similar magnitude. I therefore interpret the result as severity stratification rather than identification of dissociable cognitive profiles. The tiers were associated with clinical-unit placement (§\ref{sec:h3_results}), but the analysis did not establish added predictive value beyond the broad diagnostic grouping.



Moreover, my results inform a broader question in neuropsychology: whether variability across test batteries is primarily domain-specific or dominated by a general factor. I found that in this routinely collected cohort, the general severity axis accounted for most of the recoverable structure. Secondary domain-specific variation was present; however, I did not find that it determined the clustering. I interpret this pattern as consistent with both clinical experience, where global impairment often separates patients before they are differentiated further through narrow dissociations [e.g., \parencite{maas2017traumatic}], and the positive manifold of inter-test correlations found in neuropsychology \parencite{lezak2012neuropsychological}. Therefore, the phenotyping result I report here is more conservative than suggested by some other accounts. Specifically, when I analyzed the data under imputation-robust, sensitivity-tested conditions, overall severity emerged before specific patterns of profiles. I do not over-read this null result: domain aggregation and density clustering on a manifold dominated by a general factor have limited sensitivity to recover narrow dissociations, including an attention-specific subgroup, if such dissociations are small or orthogonal to the principal severity gradient.



My findings do not refute studies that report multiple dissociable cognitive profiles after brain injury (e.g., \parencite{garcia2020data}). Rather, they qualify their interpretations. Unlike those studies, my pipeline differed in several ways. For example, UMAP+HDBSCAN allows for processing of non-convex structures and low-density noise; the bootstrap analysis provided an evaluation of the replicability of the clusters; and each variable was standardized and winsorized in addition to being direction-aligned before forming aggregates. Each of these steps made a difference. High variance raw measures (i.e., a timed Trail making score) can dominate an unstandardized composite; non-cognitive variables can establish artificial axes; nearly empty records can disrupt density estimates. Thus apparent phenotypes may represent artifacts of preprocessing rather than true neuropsychological structure. Chapter~\ref{Chapter6} will examine this relationship further.



I regard H1 as partially supported (§\ref{sec:h1_results}). All ten imputation methods met the pre-specified silhouette and noise criteria, supporting recovery of ordered severity bands. However, overall core Jaccard ranged from 0.509 to 0.630, individual matched-cluster values ranged from 0.383 to 0.633, and the number of recovered groups varied from three to five. These results do not support sharply discrete, uniformly reproducible phenotypes. Boundary movement can occur without disrupting the underlying severity order.



Therefore, I argue that the three-tier solution developed from the primary MICE pipeline should be viewed as an interpretable discretization of a continuum rather than as evidence of natural kinds. Other clustering techniques produce between 3 and 5 tiers depending on how they are configured. Across-imputation and threshold sensitivity analyses show that cluster counts may vary while severity ordering remains consistent. The clinically useful entity is thus the continuum; the tiers represent a convenient way to describe aspects of it.

\section{Imputation Sensitivity Under a Fixed Partition}
I regard Hypothesis 2 as the major methodological contribution, but I define its scope deliberately narrowly. I discovered the primary clusters with UMAP+HDBSCAN. For the H2 analysis, I then used a separate fixed-reference diagnostic, fitting \(k\)-means once on the MICE embedding and projecting the other imputation methods into that same structure. Under that diagnostic, point assignments changed very little with regard to imputation methods: on average, the 45 method pairs yielded a mean pairwise ARI \(> 0.70\), with 100\% yielding an ARI \(> 0.50\) (§\ref{sec:h2_results}). Therefore, I interpret the easily recoverable structure as a strong severity gradient, whereas a fragile multi-dimensional profile structure is less readily recovered. This distinction likely explains the robustness while still allowing for method-specific variability \parencite{rubin1987multiple}.



Although the imputation process influenced the number of identified tiers and the quality of those tiers, I found that the location of patients along the severity signal did not appear to depend on a particular imputation strategy. I interpret the broad implication as follows: most clinical interpretations of a patient's relative position on the severity gradient are unlikely to change solely because one defensible imputation method is substituted for another. I do not take this to mean that missing data can be handled casually. Rather, I argue that researchers should treat missingness systematically and report the method clearly, while clinicians should interpret tier labels as aids to judgement rather than as rules that replace medical training or professional experience.



As depicted in Figure~\ref{fig:h2_dendrogram}, I found a strong separation between two families of imputation methods: smooth, model-based methods (EM, DAE, and VAE) that tend to smooth towards a model based conditional mean; local donor-methods (KNN) that rely on a few local donors to plug in missing values. Although many studies use classification criteria related to whether or not an imputation method uses statistical/machine learning/deep learning families, I found the more salient divide to be between smooth and local.



At the assessment level, I found that instability correlates moderately with the fraction of imputed data (\(r = 0.47\); §\ref{sec:h2_results}). As such, I argue that downstream applications should treat assessments with heavy amounts of missing data differently from assessments that contain almost no missing data. This could be achieved by reporting decreased confidence or by comparing more than one imputation strategy.



I view consensus clustering as a practical means of addressing this issue. The ten imputation strategies align their votes regarding label assignments for individual assessments, and the resulting confidence score estimates agreement among strategies. Overall confidence is high (mean confidence 0.93; 86.2\% of assessments have high confidence). However, as I discussed previously, I interpret the lower-confidence assessments as clinically meaningful rather than merely inconvenient. Most notably, these assessments represent patients who are located near tier boundaries and should prompt caution/re-assessment.



In the literature, MICE is commonly used prior to clustering neuropsychological data. MissForest \parencite{stekhoven2012missforest} is often preferred when non-linear relationships are expected. After applying Holm-Bonferroni correction across the 45 pairs of comparison (see Section~\ref{sec:robustness}), I found that 44 of the 45 pairs reject the null hypothesis that ARI \(\leq 0.50\). Only the KNN-SoftImpute pair remains at or near threshold levels. I interpret the broader result, within the fixed-reference H2 design, as clear: the location of patients along the severity signal does not depend upon a single imputation strategy. The number of density-based tiers remains more sensitive when UMAP+HDBSCAN is refitted independently.



I also provided a listwise deletion baseline in Section~\ref{sec:robustness} to illustrate why imputation is necessary. Restricting analysis to only complete cases results in a different number of clusters and silhouette (0.257) compared to the MICE solution for shared patients (ARI 0.625). Thus, while removing the 46.9\% of assessments with missing data changes recovered structure, imputation was not merely cosmetic.

\section{Clinical Unit Association}
To determine whether the severity tiers were associated with clinical-unit assignment, I tested the contingency between tier and unit. The association was statistically significant, although the effect size was small by Cohen's standards. A stratified sensitivity analysis remained significant after accounting for the broad diagnostic grouping supplied by the unit lookup. Because this grouping and detailed unit assignment originate from the same lookup, I do not treat that result as independent validation.



The association is clinically plausible, but it is not evidence of prediction or causation. In the patient-grouped multinomial sensitivity analysis, adding tier to the broad diagnostic grouping changed held-out log loss from 0.6064 to 0.6051, with a patient-bootstrap improvement interval that included zero. Accuracy was unchanged to three decimal places. I therefore interpret H3 only as evidence that tier distributions differ across units. Whether severity tier improves prospective routing or outcome prediction requires independent predictors, patient-grouped validation, and clinical outcomes collected for that purpose.

\section{Feature Engineering}
There are many ways to define cognitive features. Two broad options are variable-level features and aggregate features. Variable-level features treat each individual test separately, whereas aggregate features represent averaged or summed versions of tests that are intended to measure a shared cognitive domain. I tested whether domain-aggregate features would be better conditioned than variable-level features. In this context, better conditioned means lower multicollinearity, indexed by smaller variance inflation factors (VIFs), and a smaller condition number.



The results supported H4. I found that domain-aggregate features were much better conditioned than variable-level features even though separation between clusters increased very little. The silhouette changed only slightly (0.481 vs. 0.473); however, mean VIF decreased from 2.77 to 1.88 and the condition number decreased from 61.7 to 11.8 (see §\ref{sec:h4_results}). I therefore interpret H4 primarily as evidence that aggregation reduces redundant information and produces a more interpretable, better-conditioned representation. It does not provide evidence of materially stronger clustering. This benefit supports the clinical practice of interpreting test batteries by domain rather than test by test \parencite{lezak2012neuropsychological}.



I emphasize one critical requirement for aggregation: preprocessing must occur before aggregation. Variables must be winsorized, direction-aligned, and standardized before they can be averaged into composites. Otherwise, differences in raw scale will produce artifacts. For example, if several memory tests are averaged without standardization, the resulting composite can be dominated by whichever test has the largest range rather than by the intended memory construct. If several highly correlated tests enter into an unstandardized composite, then UMAP and HDBSCAN can exaggerate distance computation using a common cognitive dimension. Aggregation provides one feature per construct, so each domain contributes more evenly.

\section{Feature Importance and What Drives Severity}
I view the findings presented in §\ref{sec:feature_importance} as descriptive rather than explanatory. The Random Forest surrogate produced 99.4\% accuracy using the same domain scores used to generate the labels. Thus, I interpret the ranking as primarily recovering partition boundaries rather than explanatory structures underlying those partitions. The Random Forest ranked orientation, visuoperceptual abilities, and language above attention and executive function; however, I do not give this ordering strong interpretation because the domains are highly correlated. Gini importance is susceptible to variability under correlated predictors \parencite{boulesteix2012overview}.



I consider the principal component structure more reliable. One component explained 56\% of total variance in overall cognitive severity and loaded positively with similar magnitude across all six domains. I interpret this finding as evidence that severity tiers represent general severity levels rather than domain-specific dissociations or attention-specific patterns within this particular domain-aggregated pipeline. Robustness checks presented in §\ref{sec:robustness}, including subsampling and sensitivity to missingness thresholds, suggest that severity ordering remains stable, but exact tier count depends on resolution criteria and inclusion threshold. Data-driven studies examining phenotyping should report preprocessing choices clearly and include sensitivity analyses.

\section{Clinical Implications}
I interpret the current results as consistent with a three-tier cognitive severity stratification system. However, I also identify several limits to this approach. The broad clinical implication is conservative: these tiers may help summarize relative cognitive severity, but they should not be used as stand-alone clinical rules. When patients move through the rehabilitation system, clinicians will continue to rely on medical training, local service knowledge, and professional experience when making placement decisions. At present, I do not think these results justify specific recommendations for direct clinical implementation beyond careful interpretation of cognitive severity and transparent handling of missing data.



Clinicians have long relied on visual inspection of patient distributions to identify unusual cases. I interpret the current results as supporting that practice alongside algorithmic clustering. Many assessments did not fall neatly into sharply separated clusters, and I view the central tier as especially consistent with a boundary region cut from a continuum. Because the majority of patients fall along a broad severity gradient rather than within highly separated natural categories, future tools should help identify difficult-to-cluster or boundary-near assessments rather than forcing every assessment into a definitive phenotype label.



Future researchers studying large datasets with missing cognitive data may benefit from additional diagnostics for clustering difficulty. One simple candidate is the proportion of points with very low silhouette values, for example less than 0.05, reported alongside overall silhouette and stability metrics. This would not replace substantive clinical interpretation, but it could help quantify how many assessments sit near ambiguous regions of the partition.
